% --- Използвани понятия (незадължителна таблица, по образец) --
\chapter*{Използвани понятия}
\addcontentsline{toc}{chapter}{Използвани понятия}

Таблица~\ref{tab:glossary} съдържа списък на някои от използваните в тази
работа понятия, които не са изрично дефинирани в текста.

\begin{table}[H]
\centering
\renewcommand{\arraystretch}{1.3}
\begin{tabular}{@{}p{0.28\textwidth} p{0.64\textwidth}@{}}
\toprule
\textbf{Понятие} & \textbf{Описание} \\
\midrule
Динамика (\eng{velocity}) & Сила/интензитет на удара на нота, кодиран в MIDI като цяло число в диапазона 0--127. \\
MIDI & Протокол за цифрово представяне на музикални събития (ноти, време, канали, динамика). \\
E-GMD & \eng{Expanded Groove MIDI Dataset} — набор от изпълнения на електронни барабани. \\
Признак (\eng{feature}) & Числово описание, извлечено от нотата и нейния контекст, подходящо за модела. \\
Груув (\eng{groove}) & Ритмичното „усещане'' на партията, оформено от вариациите на времето и динамиката. \\
Призрачна нота (\eng{ghost note}) & Много тих удар (обикновено по малкия барабан), който запълва груува, без да е самостоятелен акцент. \\
Метрична фаза & Позицията на нотата в рамките на такт или време, изразена като число в $[0,1)$; кодира се циклично чрез $\sin/\cos$. \\
Глас (\eng{voice}) & Каноничен барабанен инструмент след категоризация на кодираните MIDI височини. \\
LightGBM & Метод за градиентно усилени дървета на решения, използван тук като базов модел за сравнение. \\
Трансформър (\eng{transformer}) & Невронна архитектура, базирана на „механизъм за внимание'' (\eng{self-attention}). \\
MDN & Мрежа със смес от гаусиани (\eng{mixture density network}) — предсказва многомодално разпределение. \\
% TODO: добави още понятия при нужда
\bottomrule
\end{tabular}
\caption{Списък на използваните понятия}
\label{tab:glossary}
\end{table}

\clearpage
